% Auto-generated by build/sm_branch_successor.py. Do not edit by hand.

\section*{SM Branch Successor Pilot}

\addcontentsline{toc}{section}{SM Branch Successor Pilot}



This generated packet is the v30 L2 branch-pilot shell for the Standard Model super-branch. It is \emph{not} a supersession of the canonical SM Branch Book. The pilot exercises the \textbf{intrinsic-structural-close} posture for the first time: SM has closed its internal structural obligations while remaining open on inherited scope from the YM and GR sub-branches.



\subsection*{Posture and inheritance structure}

Computed SM posture: \textbf{intrinsic-structural-close}. This posture encodes: \emph{(i)} conditional intrinsic structural closure (four internal obligations discharged at structural-conditional level); \emph{(ii)} inherited-scope open (all YM and GR obligations propagate to SM); \emph{(iii)} active frontier \texttt{{ob:SM.O-IDcont}} (continuum/operator identification force upgrade, the single named irreducible frontier of the SM branch).



\subsection*{Status inequality}

The governing constraint is: $\operatorname{Status}(\mathrm{SM}) \le \operatorname{Status}(\mathrm{YM}) \wedge \operatorname{Status}(\mathrm{GR}) \wedge \operatorname{Status}(\mathrm{SM}_{\mathrm{intrinsic}})$. The pilot preserves this inequality explicitly. SM cannot be promoted to conditional-cert-close or cert-close by any discharge internal to the SM packet.



\subsection*{Non-claims (strictly maintained)}

\begin{itemize}

\item No measured particle masses

\item No CKM/PMNS mixing parameters

\item No renormalization-group running derivation

\item No full empirical Standard Model reconstruction

\item No gauge-gravity unification claim

\item No unconditional CERT-CLOSE

\end{itemize}



\begin{longtable}{p{0.30\linewidth}p{0.14\linewidth}p{0.12\linewidth}p{0.15\linewidth}p{0.18\linewidth}}
\caption{SM successor-pilot claim catalog}\label{tab:sm-successor-claims}\\
\textbf{ID} & \textbf{Kind} & \textbf{Decl.} & \textbf{Comp.} & \textbf{Scope}\\
\hline\endfirsthead
\textbf{ID} & \textbf{Kind} & \textbf{Decl.} & \textbf{Comp.} & \textbf{Scope}\\
\hline\endhead
ob:SM.harmonization & obligation & C & C & scoped-KPO\\
ob:SM.matter & obligation & C & C & scoped-candidate-assignment\\
ob:SM.higgs & obligation & C & C & scoped-B0-screening\\
ob:SM.coupling-transfer & obligation & C & C & structural-seed\\
thm:SM.gauge-chain & theorem & C & C & branch\\
thm:SM.three-gen & theorem & U & U & branch\\
thm:SM.harmonization-partial & theorem & C & C & scoped-KPO\\
thm:SM.coupling-transfer-full & theorem & C & C & structural-seed\\
thm:SM.higgs-source-descended & theorem & C & C & scoped-B0-screening\\
thm:SM.18.1 & theorem & C & C & branch\\
prop:SM.intrinsic-closure & proposition & C & C & intrinsic-structural\\
prop:SM.status-reconciliation & proposition & C & C & intrinsic-structural\\
ob:SM.O-IDcont & obligation & O & O & branch\\
status:SM & status\_note & intrinsic-structural-close & intrinsic-structural-close & branch\\
\end{longtable}



\subsection*{Generated claim shells}

\begin{nfcobligation}{Harmonization Functor (O-SM.Harmonization)}
\textbf{Machine ID.} \texttt{ob:SM.harmonization}\\
\textbf{Declared.} C\quad\textbf{Computed.} C\\
\textbf{Scope.} regime=SM-super-sector; domain=SM-branch; certificate\_scope=scoped-KPO; promotion\_ceiling=C

Construct the SM harmonization functor F\_SM: it must take the pair (E\_YM, E\_GR) of sub-branch endpoint data and produce a six-component SM endpoint datum (gauge algebra, mass gap, metric class, gravitational coupling, cosmological constant, normalization unit) satisfying source-descendedness, no hidden supplementation, and the Book VII scoped-certificate condition. The KPO chain (KPO-1 through KPO-4) is the canonical candidate; extending to the full matter sector requires O-SM.Matter and O-SM.Higgs at unconditional force.

\medskip
\noindent\textit{Dependencies.} thm:YM.7, thm:GR.CP.1
\end{nfcobligation}



\begin{nfcobligation}{Matter Content (O-SM.Matter)}
\textbf{Machine ID.} \texttt{ob:SM.matter}\\
\textbf{Declared.} C\quad\textbf{Computed.} C\\
\textbf{Scope.} regime=SM-super-sector; domain=SM-branch; certificate\_scope=scoped-candidate-assignment; promotion\_ceiling=C

Assign the SM fermionic matter content (quark doublet, charged antilepton, neutrino, and their conjugates across three generations) from the NFC three-generation structure without importing particle physics primitives. The eigenblock decomposition R\textasciicircum{}9 = B\_\{+1\} + B\_\{-1\} + B\_0 provides the structural constraint; full assignment requires confirming no additional eigenmodes carry matter quantum numbers.

\medskip
\noindent\textit{Dependencies.} thm:SM.gauge-chain, thm:SM.three-gen
\end{nfcobligation}



\begin{nfcobligation}{Higgs/EWSB Screening (O-SM.Higgs)}
\textbf{Machine ID.} \texttt{ob:SM.higgs}\\
\textbf{Declared.} C\quad\textbf{Computed.} C\\
\textbf{Scope.} regime=SM-super-sector; domain=SM-branch; certificate\_scope=scoped-B0-screening; promotion\_ceiling=C

Identify the B\_0 block (the d\_\{S\_v\}=0 eigenblock, kernel of the screening projection) as the SM Higgs candidate. The screening projection must be source-descended from the NFC three-generation structure, carry the correct gauge quantum numbers under su(3)+su(2)+u(1), and produce a nonzero expectation value that breaks su(2)+u(1) to u(1)\_EM. Full ob:SM.higgs discharge requires proving the screening projection is the unique admissible Higgs mechanism at this scope.

\medskip
\noindent\textit{Dependencies.} thm:SM.gauge-chain
\end{nfcobligation}



\begin{nfcobligation}{Coupling Transfer (O-SM.CouplingTransfer)}
\textbf{Machine ID.} \texttt{ob:SM.coupling-transfer}\\
\textbf{Declared.} C\quad\textbf{Computed.} C\\
\textbf{Scope.} regime=SM-super-sector; domain=SM-branch; certificate\_scope=structural-seed; promotion\_ceiling=C

Transfer the structural gauge coupling seed from the shared archetype interface to a computable coupling constraint. The interface I (shared by YM and GR sub-branches) carries a structural coupling ratio; the seed is: the relative coupling constants are constrained by the shared interface structure. Full discharge requires deriving the numerical coupling ratios from the interface geometry, which is a deeper structural task than the present scoped discharge provides.

\medskip
\noindent\textit{Dependencies.} thm:SM.gauge-chain, thm:SM.harmonization-partial
\end{nfcobligation}



\begin{nfctheorem}{SM Gauge Algebra and Uniqueness Chain}
\textbf{Machine ID.} \texttt{thm:SM.gauge-chain}\\
\textbf{Declared.} C\quad\textbf{Computed.} C\\
\textbf{Scope.} regime=SM-super-sector; domain=SM-branch; certificate\_scope=branch; promotion\_ceiling=C; constraints=YM-Phase-A,YM-K0-7,YM-O-ENC

Starting from the YM branch endpoint E\_YM, the SM gauge algebra is identified as su(3) + su(2) + u(1) via the collar-local Lie theory chain applied to the observable super-sector. The gauge algebra is unique at this scope: no reductive compact Lie algebra other than su(3)+su(2)+u(1) admits the declared collar-local structure with the Weyl encoding constraint (O-ENC) and the K\_0=7, Phase-A hypotheses. The identification inherits full conditional force from the YM endpoint.

\medskip
\noindent\textit{Inherited constraints.} YM-Phase-A, YM-K0-7, YM-O-ENC

\medskip
\noindent\textit{Dependencies.} thm:YM.7, thm:YM.ENC, thm:YM.GLOB

\medskip
\noindent\textit{Proof stub.} Apply thm:YM.7 (YM endpoint including observable gauge algebra) and thm:YM.ENC (Weyl encoding) to the collar-local Lie sector. The collar-local Lie theory chain (Section 5 of the SM branch book) derives su(3)+su(2)+u(1) as the unique reductive compact Lie algebra fitting the declared eigenblock decomposition R\textasciicircum{}9 = B\_\{+1\}+B\_\{-1\}+B\_0 under the NFC-3Gen structure. Uniqueness follows from thm:YM.GLOB (globalization of the gauge structure): the observable gauge algebra is fully determined by the local collar data plus globalization.
\end{nfctheorem}



\begin{nfctheorem}{Three-Generation Structure}
\textbf{Machine ID.} \texttt{thm:SM.three-gen}\\
\textbf{Declared.} U\quad\textbf{Computed.} U\\
\textbf{Scope.} regime=SM-super-sector; domain=SM-branch; certificate\_scope=branch; promotion\_ceiling=U; constraints=NFC-3Gen-declared-structure

The d\_\{S\_v\}-spectral decomposition of the SM matter sector yields exactly three isomorphic copies of the matter eigenblock structure. This is unconditional: the three-generation count follows from the rank of the structural eigenspace decomposition and is independent of the conditional YM/GR hypotheses. It is the structural origin of three matter generations in the NFC framework.

\medskip
\noindent\textit{Inherited constraints.} NFC-3Gen-declared-structure

\medskip
\noindent\textit{Proof stub.} The eigenspace decomposition of the SM matter sector under d\_\{S\_v\} produces eigenblocks B\_\{+1\}, B\_\{-1\}, B\_0 each of multiplicity three (from the NFC-3Gen structure established in thm:SM.gauge-chain). The three copies are isomorphic by the symmetry of the eigenblock decomposition. No conditional YM or GR hypothesis is needed: the three-generation count is a consequence of the declared collar structure, which is unconditional.
\end{nfctheorem}



\begin{nfctheorem}{Conditional Partial Harmonization Functor (KPO-1 through KPO-4)}
\textbf{Machine ID.} \texttt{thm:SM.harmonization-partial}\\
\textbf{Declared.} C\quad\textbf{Computed.} C\\
\textbf{Scope.} regime=SM-super-sector; domain=SM-branch; certificate\_scope=scoped-KPO; promotion\_ceiling=C; constraints=YM-Phase-A,YM-K0-7,YM-O-ENC,GR-KPO-3,GR-NFC-5-6

The KPO chain stages KPO-1 through KPO-4 define a candidate harmonization functor F\textasciicircum{}KPO\_SM: F\textasciicircum{}KPO\_SM(E\_YM, E\_GR) = (g\_obs, m\textasciicircum{}2, [g\_\{mu nu\}], kappa, Lambda, beta), where each component is inherited from declared sub-branch data without new target-side choices. The functor satisfies source-descendedness and no-hidden-supplementation conditionally. Remaining: extending to the full matter sector (O-SM.Matter, O-SM.Higgs) and formal POT-category packaging as an explicit morphism. O-SM.Harmonization is therefore partially discharged at this scope.

\medskip
\noindent\textit{Inherited constraints.} YM-Phase-A, YM-K0-7, YM-O-ENC, GR-KPO-3, GR-NFC-5-6

\medskip
\noindent\textit{Dependencies.} thm:YM.7, thm:GR.CP.1, thm:SM.gauge-chain

\medskip
\noindent\textit{Discharge edges.} ob:SM.harmonization (partial)

\medskip
\noindent\textit{Proof stub.} Source-descendedness: each component of F\textasciicircum{}KPO\_SM is inherited from E\_YM or E\_GR via the declared KPO stages. No new target-side structure is introduced. No-hidden-supplementation: the KPO chain uses only declared canonical primitives at each stage; the six-component output is fully determined by (E\_YM, E\_GR). Conditional force: the functor carries [C] status bounded by the conjunction of YM hypotheses (Phase-A, K\_0=7, O-ENC, O-ID, O-GLOB, O-CLU) and GR hypotheses (KPO-3, NFC-5/6, curvature-subcriticality). Partial discharge mode: the matter sector (O-SM.Matter, O-SM.Higgs) is not yet in E\_SM.
\end{nfctheorem}



\begin{nfctheorem}{Structural Coupling Transfer Theorem}
\textbf{Machine ID.} \texttt{thm:SM.coupling-transfer-full}\\
\textbf{Declared.} C\quad\textbf{Computed.} C\\
\textbf{Scope.} regime=SM-super-sector; domain=SM-branch; certificate\_scope=structural-seed; promotion\_ceiling=C

The shared archetype interface I (used by both YM and GR sub-branches) carries a structural coupling seed: the relative coupling constants are constrained by the interface geometry without free parameters at this scope. The structural 6:2:1 coupling ratio is a consequence of the interface geometry. This is the structural seed of gauge coupling unification: it constrains the coupling constants without deriving the numerical values from first principles. The scoped discharge covers the structural-seed level only; full coupling constant derivation remains at deeper scope.

\medskip
\noindent\textit{Dependencies.} thm:SM.harmonization-partial, thm:YM.7, thm:GR.CP.1

\medskip
\noindent\textit{Discharge edges.} ob:SM.coupling-transfer (scoped)

\medskip
\noindent\textit{Proof stub.} The interface I is shared by the YM archetype super-sector (SA\_1) and the GR archetype super-sector (SA\_2). Structural coupling: the relative coupling of the two interfaces is constrained by the shared interface geometry (prop:SM-interface-coupling in the canonical SM branch). The 6:2:1 ratio follows from the eigenblock dimensions. Scoped: this is a structural constraint, not a numerical derivation of the measured Standard Model coupling constants. The full coupling constant derivation would require O-IDcont and deeper structural analysis.
\end{nfctheorem}



\begin{nfctheorem}{B0 Screening Projection Source-Descendedness}
\textbf{Machine ID.} \texttt{thm:SM.higgs-source-descended}\\
\textbf{Declared.} C\quad\textbf{Computed.} C\\
\textbf{Scope.} regime=SM-super-sector; domain=SM-branch; certificate\_scope=scoped-B0-screening; promotion\_ceiling=C

The B\_0 screening projection as Higgs mechanism is source-descended: the screening projection onto B\_0 = ker(d\_\{S\_v\}) is uniquely determined by the declared NFC-3Gen structure without external choice. B\_0 carries the correct gauge quantum numbers under su(3)+su(2)+u(1) (transforming as (1,1,0) i.e. a gauge singlet of the strong and weak interactions). When B\_0 carries a nonzero expectation value, the residual gauge symmetry is u(1)\_EM. Source-descendedness: B\_0 = ker(d\_\{S\_v\}) is uniquely determined by the declared eigenblock decomposition.

\medskip
\noindent\textit{Dependencies.} thm:SM.gauge-chain, thm:SM.three-gen

\medskip
\noindent\textit{Discharge edges.} ob:SM.higgs (scoped)

\medskip
\noindent\textit{Proof stub.} B\_0 = ker(d\_\{S\_v\}) is uniquely determined by the NFC-3Gen structure (thm:SM.three-gen): it is the unique zero-eigenvalue block of the spectral decomposition. No external definition is needed. Gauge quantum numbers: the gauge algebra identification (thm:SM.gauge-chain) assigns B\_0 to the singlet representation under su(3)+su(2). EWSB: the B\_0 expectation value breaks the su(2)+u(1) factor to u(1)\_EM via the standard Higgs mechanism analog; the screening projection is the unique admissible projection that preserves su(3)+u(1)\_EM. Scoped: full ob:SM.higgs requires proving the screening projection is the unique admissible mechanism, which needs deeper structural analysis.
\end{nfctheorem}



\begin{nfctheorem}{Conditional SM Intrinsic Closure}
\textbf{Machine ID.} \texttt{thm:SM.18.1}\\
\textbf{Declared.} C\quad\textbf{Computed.} C\\
\textbf{Scope.} regime=SM-super-sector; domain=SM-branch; certificate\_scope=branch; promotion\_ceiling=C

The SM super-branch is conditionally intrinsic-structurally closed, subject to inherited-scope and force-upgrade limits.

\medskip
\noindent\textit{Dependencies.} thm:YM.7, thm:GR.CP.1, thm:SM.harmonization-partial, thm:SM.higgs-source-descended, thm:SM.coupling-transfer-full
\end{nfctheorem}



\begin{nfcproposition}{Conditional SM Intrinsic Structural Closure}
\textbf{Machine ID.} \texttt{prop:SM.intrinsic-closure}\\
\textbf{Declared.} C\quad\textbf{Computed.} C\\
\textbf{Scope.} regime=SM-super-sector; domain=SM-branch; certificate\_scope=intrinsic-structural; promotion\_ceiling=C; constraints=YM-all-open-obligations,GR-curvature-subcriticality

The SM super-branch has conditional intrinsic structural closure: its internal anti-smuggling obligations for gauge algebra, gauge-template uniqueness, YM-GR interface compatibility, harmonization functor (KPO scope), matter datum (candidate assignment), Higgs/EWSB screening (B\_0 scope), and structural coupling transfer are discharged at conditional force.

This proposition does NOT promote SM to unconditional CERT-CLOSE. SM: (1) inherits all unresolved obligations of the YM and GR sub-branches; (2) does not claim low-energy phenomenological matching, RG-running derivation, measured particle masses, CKM/PMNS data, or full empirical Standard Model reconstruction; (3) does not claim unification of gauge and gravitational forces.

The status inequality: Status(SM) <= Status(YM) AND Status(GR) AND Status(SM\_intrinsic).

\medskip
\noindent\textit{Inherited constraints.} YM-all-open-obligations, GR-curvature-subcriticality

\medskip
\noindent\textit{Dependencies.} thm:SM.harmonization-partial, thm:SM.higgs-source-descended, thm:SM.coupling-transfer-full, thm:SM.18.1

\medskip
\noindent\textit{Discharge edges.} ob:SM.harmonization (scoped), ob:SM.matter (scoped), ob:SM.higgs (scoped), ob:SM.coupling-transfer (scoped)

\medskip
\noindent\textit{Proof stub.} The four SM-internal obligations are conditionally discharged: O-SM.Harmonization by thm:SM.harmonization-partial (partial, KPO scope); O-SM.Matter by candidate assignment at structural level; O-SM.Higgs by thm:SM.higgs-source-descended (B\_0 scope); O-SM.CouplingTransfer by thm:SM.coupling-transfer-full (structural-seed scope). All discharges are scoped: they do not reach unconditional force. The inherited-scope clause follows from the super-branch governance rule: the SM super-branch harmonizes YM and GR endpoint data and cannot be stronger than the conjunction of its parent-branch assumptions.
\end{nfcproposition}



\begin{nfcproposition}{SM Status Reconciliation}
\textbf{Machine ID.} \texttt{prop:SM.status-reconciliation}\\
\textbf{Declared.} C\quad\textbf{Computed.} C\\
\textbf{Scope.} regime=SM-super-sector; domain=SM-branch; certificate\_scope=intrinsic-structural; promotion\_ceiling=C; constraints=YM-inherited-scope,GR-inherited-scope

SM-New supersedes SM-Old (CERT-PROJ) only with respect to the internal SM obligation ledger. The corrected status is: SM: conditionally intrinsic-structural closed, inherited-scope open.

This is strictly stronger than raw CERT-PROJ (the four internal obligations are conditionally discharged), but strictly weaker than unconditional or full empirical CERT-CLOSE.

The following non-claims remain fully binding: (1) inherited YM/GR obligations and scope limits; (2) no unconditional CERT-CLOSE; (3) no phenomenological Standard Model completion (measured masses, CKM/PMNS, RG-running, full empirical reconstruction); (4) no gauge-gravity unification claim.

The next highest-leverage obligation is the extension of the harmonization functor to the full matter sector via O-IDcont force upgrade: F\textasciicircum{}\{KPO+M\}\_SM, requiring O-ID\textasciicircum{}cont at unconditional force from the YM branch.

\medskip
\noindent\textit{Inherited constraints.} YM-inherited-scope, GR-inherited-scope

\medskip
\noindent\textit{Dependencies.} prop:SM.intrinsic-closure, ob:SM.O-IDcont

\medskip
\noindent\textit{Proof stub.} prop:SM.intrinsic-closure gives the internal structural closure. The inherited-scope and non-claims clauses follow from the super-branch governance rule: the SM super-branch harmonizes YM and GR endpoint data and cannot be stronger than the conjunction of its parent-branch assumptions. The non-claims list is carried forward from the SM branch book boundary declaration. The status ordering Status(SM) <= Status(YM) AND Status(GR) AND Status(SM\_intrinsic) is an immediate consequence.
\end{nfcproposition}



\begin{nfcobligation}{Continuum/Operator Identification Force Upgrade}
\textbf{Machine ID.} \texttt{ob:SM.O-IDcont}\\
\textbf{Declared.} O\quad\textbf{Computed.} O\\
\textbf{Scope.} regime=SM-super-sector; domain=SM-branch; certificate\_scope=branch

The continuum/operator identification force upgrade: upgrading the O-ID\textasciicircum{}cont identification from its current conditional/structural scope to unconditional force. Without this, the SM harmonization functor extension to the full matter sector (F\textasciicircum{}\{KPO+M\}\_SM) cannot reach unconditional force, and SM intrinsic structural closure cannot be promoted to unconditional CERT-CLOSE. This obligation directly inherits from the YM branch ob:O-YM-O-IDcont and is the principal force-upgrade frontier of the SM branch.

\medskip
\noindent\textit{Dependencies.} prop:SM.intrinsic-closure
\end{nfcobligation}



\begin{nfcstatusnote}{SM Branch Posture}
\textbf{Machine ID.} \texttt{status:SM}\\
\textbf{Declared.} intrinsic-structural-close\quad\textbf{Computed.} intrinsic-structural-close\\
\textbf{Scope.} regime=SM-super-sector; domain=SM-branch; certificate\_scope=branch

SM branch posture: intrinsic-structural-close (conditionally intrinsic-structural closed, inherited-scope open).

This is the first pilot to exercise the intrinsic-structural-close posture. It is distinct from conditional-cert-close: the SM branch is not simply conditionally approaching a known endpoint, but has closed its internal structural obligations while remaining open on inherited scope from YM and GR.

Active frontier: ob:SM.O-IDcont remains open. The continuum/operator identification force upgrade is the single named irreducible frontier of the SM branch. Without O-IDcont at unconditional force, the SM branch cannot achieve unconditional closure even if all internal obligations are discharged.

Non-claims strictly maintained: no measured particle masses, no CKM/PMNS, no RG-running, no full empirical Standard Model reconstruction, no gauge-gravity unification.

The status inequality Status(SM) <= Status(YM) AND Status(GR) AND Status(SM\_intrinsic) is the governing constraint. The pilot preserves this inequality explicitly.

\medskip
\noindent\textit{Dependencies.} prop:SM.intrinsic-closure, prop:SM.status-reconciliation, ob:SM.O-IDcont, thm:V.UBLT, thm:VII.6.4
\end{nfcstatusnote}


